\documentclass[11pt]{article}
\usepackage[T1]{fontenc}
\usepackage[utf8]{inputenc}
\usepackage{lmodern,microtype}
\usepackage[margin=1in]{geometry}
\usepackage{amsmath,amssymb,amsthm,mathtools}
\usepackage{booktabs,longtable,array,enumitem}
\usepackage{xcolor}
\usepackage{hyperref}
\usepackage{listings}
\hypersetup{colorlinks=true,linkcolor=blue!45!black,urlcolor=blue!45!black,citecolor=blue!45!black,pdftitle={A prefix presentation of intuitionistic primitive recursive arithmetic},pdfauthor={Research prepared for Vladimir Reshetnikov}}
\lstset{basicstyle=\small\ttfamily,breaklines=true,columns=fullflexible,keepspaces=true,frame=single,rulecolor=\color{black!20}}
\setlist{itemsep=3pt,topsep=5pt}
\newtheorem{theorem}{Theorem}[section]
\newtheorem{lemma}[theorem]{Lemma}
\newtheorem{proposition}[theorem]{Proposition}
\newtheorem{corollary}[theorem]{Corollary}
\theoremstyle{definition}\newtheorem{definition}[theorem]{Definition}
\theoremstyle{remark}\newtheorem{remark}[theorem]{Remark}
\newcommand{\I}{\mathsf I}
\newcommand{\D}{\mathsf D}
\newcommand{\Q}{\mathsf Q}
\newcommand{\F}{\mathcal F}
\newcommand{\Hcal}{\mathcal H}
\newcommand{\ipra}{\mathrm{iPRA}}
\newcommand{\Teq}{T_{\rm eq}}
\newcommand{\Eq}[3]{\langle #1;\,#2=#3\rangle}
\newcommand{\Clause}[2]{\langle #1;\,y.#2\rangle}
\newcommand{\one}{S0}
\newcommand{\code}[1]{\texttt{#1}}
\newcommand{\srcbase}{https://github.com/VladimirReshetnikov/ProveIt/tree/17898ea7e6a1300dfdcad71040d5daeed67916d4/Logic/PrimitiveRecursiveArithmetic/ConciseIPRA}
\title{A prefix presentation of intuitionistic\protect\\ primitive recursive arithmetic\large\protect\\ Four rule families, one recursive object syntax,\protect\\ and an antecedent-free alternative}
\author{Research prepared for Vladimir Reshetnikov}
\date{November 11, 2026}
\begin{document}
\maketitle
\begin{abstract}
The current ConciseIPRA development already has four instruction families
with a fixed reflexivity clause. We retain its arithmetic and intuitionistic
logical strength and simplify its presentation: every native Eq/$D$ formula
is a chain of paired binders ending in one equation, so it can be represented
by a natural number and two terms. Only the numerical term syntax remains
recursive. Substitution in a formula becomes substitution in its two terms
at a depth determined by that number. We give the complete calculus, an
executable checker, explicit translations to and from the current calculus,
and the resulting whole-domain bi-interpretation with iPRA.

A separate mathematical result removes the antecedent field of every
context clause. An asymmetric iterator gate represents ordinary assumptions
as homogeneous conditional closures and has a literally reflexive zero
instance. This alternative is proved equivalent, but its measured checker
is larger than the corresponding prefix checker with paired clauses.
The package distinguishes these results from a claim of global minimality:
it does not beat the repository's most compressed existing implementations.
\end{abstract}

\section{Result, comparison target, and limits}

The inspected baseline is commit
\begin{center}\small\ttfamily
17898ea7e6a1300dfdcad71040d5daeed67916d4
\end{center}
of \emph{ProveIt}, under
\code{Logic/PrimitiveRecursiveArithmetic/ConciseIPRA}~\cite{repo}.
All comparisons in this report refer to that pinned snapshot.
The baseline already contains the clause calculi \code{CPrf}, \code{KPrf},
and \code{VPrf}. In particular, putting reflexivity into a fixed clause,
using two-case computational equality, and using four instruction families
are inherited results~\cite{clausal,uniform,eval}. They are not claimed as
new here.

The primary presentation $\F$ makes three precise choices:
\begin{enumerate}
\item A formula is $\Eq{k}{s}{t}$: a binder count and a pair of numerical
terms. There is no recursive formula constructor in the certificate grammar.
\item A clause has one antecedent and one unary consequent template. A
reflexive-leaf antecedent may be discharged without a supplied proof.
This includes the fixed reflexivity clause and ordinary assumptions.
\item There are four proof constructors: Use, Eval, Elim, and Intro.
Elim synthesizes its conclusion and checks that its local witness has not
escaped.
\end{enumerate}
The first choice is the main representational reduction. The second and
third make the complete interface uniform. Section~\ref{sec:measurements}
separates the complete comparison from the comparison in which the control
uses the same proof interface.

The arithmetic basis itself is unchanged: a scalar iterator has a seed
and a unary update, and the restricted induction mechanism is the inherited
witness-directed introduction. The new claim is a smaller explicit
presentation and implementation in the reported comparison, not a new
strength theorem about arithmetic or a minimum over all programming
languages and proof encodings. The existing compact checkers remain
smaller in absolute source size. The additional homogeneous-context theorem
is a mathematical simplification whose implementation cost is reported
rather than omitted.

\section{The complete prefix calculus}
\label{sec:calculus}

\subsection{Numerical terms and formula prefixes}

There is one numerical sort. Terms are
\[
 t ::= x\mid 0\mid St\mid\I(n,a;u.h),
\]
where only $u$ is bound in $h$. A term may contain any finite nesting of
iterators; there are no function variables or universal evaluator.

A formula is the finite record $\Eq{k}{s}{t}$, with $k\in\mathbb N$.
It binds $k$ ordered pairs of numerical variables in $s,t$. Its expansion is
\[
 \operatorname{unpack}\Eq{k}{s}{t}
   =\D_{x_1,y_1}\cdots\D_{x_k,y_k}(s=t).
 \tag{1}\label{eq:unpack}
\]
Thus $\Eq{0}{s}{t}$ is an equation. The names in (1) are explanatory;
the checker uses bound indices. The binder count is syntactic data, not an
object-level numerical variable quantified by the theory.

For readability, $\D_{x,y}B$ below denotes adding one pair to a prefix.
It is an operation on the record, not an additional recursive syntax.
$B(a,b)$ denotes capture-free opening of its two designated external slots.
A unary template $B(y)$ has one designated external slot. All other free
parameters remain fixed.

The intended ordinary reading of the prefix is recursively specified by
\[
 (s=t)^\circ=(s=t),\qquad
 (\D_{x,y}B)^\circ
  =\exists x\bigl(B^\circ(x,0)\to\forall y\,B^\circ(x,y)\bigr).
 \tag{2}\label{eq:ordinary-D}
\]
This recursion defines a translation of each finite syntactic object.
It does not assert the existence of an internal truth predicate.
The designated contradiction is $\Eq{0}{0}{S0}$.

\subsection{Contexts and their one fixed clause}

A context is a finite reusable list of clauses $[A;y.B(y)]$. Its ordinary
reading is $A^\circ\to\forall y B^\circ(y)$; $y$ is not free in $A$.
Each proof begins with the clause
\[
 \rho=[\,0=0;\ y.y=y\,].
 \tag{3}\label{eq:rho}
\]
In the closed-certificate protocol this is the sole initial entry, and
the list grows only through Intro and Elim below. For mathematical
statements about sequents, additional well-scoped clauses may be retained
as external assumptions; the delivered verifier checks closed certificates.
An ordinary local
assumption $A$ is the constant clause $[0=0;v.A]$, where $v$ is unused.
There is no separate hypothesis rule.

Call $\Eq{k}{r}{r}$ \emph{reflexive at its leaf}. The two term trees must
be literally identical. No numerical evaluation, provability test, or
normalization is used in this test. Section~\ref{sec:reflexive}
proves that every such formula is a theorem, including $k>0$.

\subsection{Four primitive rule families}

Write $\Gamma\vdash_\F A$ for derivability. Contexts always include $\rho$.
The entire primitive inventory is as follows.

\paragraph{Use.}
For any retained clause $[A;y.B]\in\Gamma$ and any term $t$,
\[
 \frac{\Gamma\vdash_\F A}{\Gamma\vdash_\F B(t)}.
 \tag{Use}
\]
The minor is omitted when $A$ is reflexive at its leaf. Otherwise one
minor derivation is required. Clause use does not consume the clause.

\paragraph{Eval.}
From two equations,
\[
 \frac{\Gamma\vdash_\F s=t\qquad\Gamma\vdash_\F b=r}
      {\Gamma\vdash_\F r=\Phi(s,t,b)},
 \tag{Eval}
\]
where $\Phi$ is defined in exactly two cases:
\[
 \Phi(s,t,b)=
 \begin{cases}
 t,&b\equiv\I(0,s;u.h),\\[2pt]
 h[t/u],&b\equiv\I(Sc,a;u.h)\text{ and }s\equiv\I(c,a;u.h).
 \end{cases}
 \tag{4}\label{eq:phi}
\]
All displayed matches are literal syntactic matches, with the usual
capture-free interpretation of a binder. There is no extra equality-join
case. The two lines of (4) are the inherited computational equality
interface~\cite{eval,clausal}.

\paragraph{Elim.}
\[
 \frac{\Gamma\vdash_\F\D_{x,y}B(x,y)\qquad
       \Gamma,[B(b,0);y.B(b,y)]\vdash_\F A}
      {\Gamma\vdash_\F A}.
 \tag{Elim}
\]
The numerical parameter $b$ is fresh for $\Gamma$, the conclusion $A$,
and the fixed parameters of the eliminated formula. The checker enforces
this by assigning a new free level and validating the branch result in
the outer scope, including occurrences in unused iterator bodies.

\paragraph{Intro.}
For $z$ fresh for $\Gamma$, $t$, and the conclusion's fixed parameters,
\[
 \frac{\Gamma,[0=0;v.B(t,\eta)]\vdash_\F B(t,Sz)}
      {\Gamma\vdash_\F\D_{x,y}B(x,y)},
 \qquad
 \eta=\begin{cases}
 z,&B\text{ is an equation and }t\equiv0,\\
 0,&\text{otherwise}.
 \end{cases}
 \tag{Intro}\label{eq:intro}
\]
Here $v$ is unused. In the prefix representation the test ``$B$ is an
equation'' means that the conclusion has binder count exactly $1$.
This is the established single witness-directed introduction; its case at
a witness that is not literally zero supplies numerical cases, and its
atomic zero-witness case supplies equation induction~\cite{repo,witness}.

There are no other primitive logical axioms, equality axioms, computation
axioms, induction schemas, structural rules, or inference rules in $\F$.
The fixed clause (3), the two syntactic cases of Eval, and the stated side
conditions are part of this complete specification.

\section{Why the reduced interface is sufficient}

\subsection{Reflexive leaves and ordinary assumptions}
\label{sec:reflexive}

\begin{lemma}[Reflexivity]
For every well-scoped term $t$, $\Gamma\vdash_\F t=t$.
\end{lemma}
\begin{proof}
Use $\rho$ at $t$. Its antecedent is the literally reflexive equation
$0=0$, so the minor is omitted.
\end{proof}

\begin{lemma}[Reflexive prefixes]
Every well-scoped formula $\Eq{k}{r}{r}$ is derivable in every context.
The same statement holds in the inherited \code{CPrf} and \code{VPrf}
calculi after unpacking. The derivations can be constructed effectively
from the formula, independently in each calculus.
\end{lemma}
\begin{proof}
Induct on $k$. The case $k=0$ is the previous lemma in $\F$, primitive
Refl in \code{CPrf}, and use of the unconditional fixed reflexivity clause
in \code{VPrf}. If $k>0$, introduce
the outer pair with witness $S0$. This is case mode regardless of the
matrix's root. The target at the fresh successor has binder count $k-1$,
and its two leaf terms are still identical after opening the witness and
counter slots. Apply the induction hypothesis to this target and ignore
the added clause. Close Intro. Consequently the optional minor in Use
only abbreviates a derivation already available in the system.
\end{proof}

Use of the constant clause $[0=0;v.A]$ gives $A$ at any target, again
without a minor. Thus the ordinary-assumption interface is available with
no separate hypothesis constructor.

\subsection{Computation and equality}

Both iterator computation equations are Eval on two reflexivities:
\[
 \I(0,a;u.h)=a,\qquad
 \I(Sn,a;u.h)=h[\I(n,a;u.h)/u].
 \tag{5}\label{eq:compute}
\]
To see that no equality principle has disappeared, let
$F(s)=\I(0,s;u.u)$. From $s=r$ and reflexivity of $F(s)$, the zero case
of Eval gives $F(s)=r$. Applying the same zero case to $s=t$ and this
last equation gives
\[
 s=t,\ s=r\ \vdash_\F\ r=t.
 \tag{6}\label{eq:join}
\]
Taking one premise reflexive yields symmetry, and symmetry plus (6)
yields transitivity.

For an arbitrary term template $h(u)$ and a premise $s=t$, the zero
case gives $\I(0,s;u.h)=t$. The successor case, with reflexivity of
$\I(S0,s;u.h)$, gives $\I(S0,s;u.h)=h(t)$. Independently, zero Eval
on $s=s$ and reflexivity of $\I(0,s;u.h)$ gives
$\I(0,s;u.h)=s$. Successor Eval with reflexivity of $\I(S0,s;u.h)$
then gives $\I(S0,s;u.h)=h(s)$. Joining these two equations yields
$h(s)=h(t)$, without first rewriting under $h$. This supplies arbitrary numerical congruence
and the inherited transport rule.

Full formula transport then follows by induction on the number of
$D$-blocks: eliminate the source paired formula, introduce the target at
the fresh witness, convert its local base, use the retained clause at the
successor, and convert back. The fresh witness forces case mode. All
matrix conversions are at smaller binder count.

\subsection{Induction is present and remains restricted}

Let $E(y)$ be a literal equation. Given proofs of $E(0)$ and of
$E(Sz)$ from the local assumption $E(z)$, Intro at the unused witness
$0$ gives $\D_{x,y}E(y)$. Elim exports $[E(0);y.E(y)]$; Use with the
base proof gives $E(t)$ for any term $t$. This is contextual equation
induction. It applies to equations, not to arbitrary formulas.

The inherited witness-rule equivalence also recovers the general case
introduction. Its exceptional atomic case at witness $0$ can use the
syntactically nonzero but provably zero term $\I(0,0;u.0)$ and the derived
transport. This explains why the literal witness test is harmless, while
also explaining why substitution admissibility is proved through the
equivalence and cannot be read off from an invariant witness test alone.

\section{Exact correspondence with the current calculus}
\label{sec:equivalence}

\begin{lemma}[Unique prefix decomposition]
Every formula in the current Eq/$D$ grammar has a unique representation
$\Eq{k}{s}{t}$, up to the ordinary representation of bound variables.
Packing and unpacking are inverse operations on these finite trees.
\end{lemma}
\begin{proof}
An equation supplies $k=0$. At a $D$ node there is exactly one formula
child; strip that node and add one to the child's count. The grammar has
no branching formula constructor. Repeatedly stripping $D$ therefore
reaches exactly one equality leaf. Reconstruction adds exactly the
stripped nodes in their original order.
\end{proof}

With de Bruijn indices starting at $1$, the two outer slots of the matrix
of a rank-$k+1$ formula have indices $2k+2$ and $2k+1$ at its equality
leaf. Its $k$ inner pairs occupy the smaller positive indices. All concrete substitution values are closed with respect to exterior
positive indices; their free numerical parameters have negative levels.
Opening
the outer witness substitutes at $2k+2$; after that, opening the outer
counter substitutes at $2k+1$. An additional iterator binder increments
the index being sought by one. This gives the following commuting law.

\begin{lemma}[Opening commutes with packing]
Packing a capture-free opening of a nested formula equals opening its
packed record by the preceding index rule. The same holds under every
iterator binder and for simultaneous choices of the external slots.
\end{lemma}
\begin{proof}
Strip the formula prefix. Each stripped pair adds exactly two to the
substitution depth. At the equality leaf, induct on its two numerical
terms. A variable is replaced exactly at the indicated index; larger
indices decrease by one when that slot is removed. Successor adds no
binder, and the update body of an iterator adds one. These are exactly
the clauses of capture-free substitution on the nested representation.
\end{proof}

\begin{theorem}[Equivalence of finite derivations]
The prefix calculus $\F$ and the Eq/$D$ (that is, $Q$-free) part of the
pinned native clause calculus derive
the same sequents after packing or unpacking formulas and clauses.
The translations of derivations are effective. Both formula round trips
are syntactic identities, not merely agreement in the standard model.
\end{theorem}
\begin{proof}
First compare with \code{CPrf}, whose contexts already consist of pairs
$[A;y.B]$ and whose ordinary assumptions are $[0=0;v.A]$.

From \code{CPrf} to $\F$, pack all formulas and add the fixed clause
$\rho$ to the context.  Primitive reflexivity becomes a zero-premise
Use of $\rho$.  The two Eval patterns, ordinary Use, Elim and the
witness-directed Intro otherwise translate directly.  Opening commutes
with packing, and the sets of free numerical parameters are unchanged,
so all substitutions and eigenvariable conditions agree.  An induction
on the source derivation preserves the extra fixed clause under every
scope extension.

For the reverse direction, carry the context invariant that every
retained clause is either the distinguished $\rho$ or the packing of a
\code{CPrf} context clause.  Use of the distinguished $\rho$ becomes
\code{CPrf} reflexivity.  Any other Use carries its translated
antecedent derivation; if that derivation was omitted, supply the
independently proved \code{CPrf} reflexive-prefix theorem instead.
Then apply the ordinary \code{CPrf} Use rule.  Eval, Elim and Intro
translate directly and preserve this invariant.  Thus no unproved
clause-cut or context-extension principle is required.

The pinned theorem \code{vprf\_iff\_cprf} gives the corresponding
four-family \code{VPrf} formulation.  Its empty antecedent is represented
by $0=0$, and its singleton antecedent remains the same formula.
The above proof explains exactly how the fixed reflexivity clause and
all skipped minors are handled; it does not assume identical certificate
grammars.  Synthesizing Elim's conclusion instead of storing another
copy preserves its outer-scope side condition.  Finally, the exact
native equivalences of \code{CPrf}/\code{VPrf} are the pinned inherited
results cited in the report.
\end{proof}

\section{Ordinary logic and the bi-interpretation with iPRA}
\label{sec:biinterpretation}

The preceding theorem is an equivalence of complete proof presentations.
Here iPRA means intuitionistic first-order logic with equality, symbols
for every finite primitive-recursive description and their defining
equations, successor separation and injectivity, and contextual induction
for quantifier-free formulas. We specify the arithmetic interpretation as well, so that the conclusion
is stronger than equiconsistency or standard-model representability.

\subsection{All ordinary logical operations are available}

The logical macros below are inherited from the native construction.
They can be packed into prefix records after expansion. Put
$C(v,a,b)=\I(v,a;u.b)$ with $u$ unused. Define
\[
 G_v(s=t)=(C(v,0,s)=C(v,0,t)),\qquad G_v(\D B)=\D(G_v B).
\]
The native calculus proves $G_0 A$ and
$G_{Sv}A\dashv\vdash A$. To select between formulas of different
depths, pad the shallower one with
$A\dashv\vdash\D_{u,y}G_yA$, using an unused witness $u$.
After their depths are aligned, define the selector at the leaf by
\[
 M_v(s=t,a=b)=(C(v,s,a)=C(v,t,b)),
\]
and pass it through aligned $D$-blocks. Its zero and positive endpoints
are $A$ and $B$, respectively. Let
$\delta(v)=C(v,1,0)$ and $p(z)=\I(z,0;u.\delta(u))$.
With all displayed dummy variables fresh, the macros are
\[
\begin{aligned}
 A\Rightarrow_* B&=\D_{u,z}M_z(A,B),\\
 A\wedge_*B&=\D_{u,z}G_zM_{p(z)}(A,B),\\
 \exists_*x\,A(x)&=\D_{x,y}G_yA(x),\\
 A\vee_*B&=\exists_*x\,M_x(A,B),\\
 \forall_*x\,A(x)&=A(0)\wedge_*\D_{u,z}G_zA(z),\\
 \bot_*&=(0=S0).
\end{aligned}
\tag{7}\label{eq:macros}
\]
The gate, selector, and paired introduction/elimination prove the usual
intuitionistic introduction and elimination rules. Numerical cases derive
$x=0\vee_*\exists_*v\,x=Sv$. Under $0=S0$, selector computation and
transport give $s=C(0,s,t)=C(S0,s,t)=t$, and induction on formula depth
gives explosion for every native formula. The inherited detailed proof is
included as \code{reference/inherited\_choice\_basis.tex}; the corresponding
Lean results are identified in \cite{repo,native-bi}.

\subsection{The ordinary scalar realization}

Let $\Teq$ be intuitionistic first-order logic with actual equality over
the scalar iterator language, with equations (5), contextual induction
for literal equations, and the sentences
\[
 \forall x\,(x=0\vee\exists y\,x=Sy),\qquad
 (0=S0)\to\bot.
 \tag{8}\label{eq:coverage}
\]
These are the ordinary counterpart of the derived native interface.
The explicit falsity bridge is essential in this ordinary description:
positive equations alone would admit a one-element structure.
The prefix-to-ordinary map is (1)--(2), and an ordinary formula is mapped
back by (7) followed by packing. Native completeness and both formula
round trips are proved in the pinned development. Packing adds an exact
syntactic isomorphism, so those results transfer to $\F$.

\subsection{The numerical interpretations}

Interpret $\I(n,a;u.h)$ in iPRA as the primitive recursion with initial
value $a$ and update $h$, ignoring the recursion counter in that update.
Each fixed nested term gives a finite primitive-recursive description.
The interpretation domain is the entire numerical domain, and its
equality is actual equality.

Conversely, the scalar theory first derives a predecessor $P$ with
$P0=0$ and $P(Sn)=n$, then a numerical test reflecting actual equality,
then quantifier-free induction, and then pairing and projections.
This order avoids assuming pairing or an equality detector in the
predecessor argument. The inherited detailed proofs are supplied in the
reference directory~\cite{bridge}.
If $\pi,L,H$ satisfy
$L(\pi(a,b))=a$ and $H(\pi(a,b))=b$, a fixed primitive recursion
\[
 f(0,\vec p)=g(\vec p),\qquad
 f(Sn,\vec p)=h(n,f(n,\vec p),\vec p)
\]
is compiled to
\[
\begin{aligned}
 V(n,\vec p)&=\I\bigl(n,\pi(0,g^*(\vec p));\,
 u.\pi(SL(u),h^*(L(u),H(u),\vec p))\bigr),\\
 f^*(n,\vec p)&=H(V(n,\vec p)).
\end{aligned}
\tag{9}\label{eq:compiler}
\]
Equation induction proves $L(V(n,\vec p))=n$, and the two projection
identities prove the required recursion equations. Pairing stores a
finite state; it does not replace the numerical universe by a range
of codes.

\begin{theorem}[Whole-domain bi-interpretability and equiconsistency]
The ordinary realization of $\F$ is bi-interpretable with full
first-order intuitionistic primitive recursive arithmetic. Both numerical
domains are whole, both equality relations are actual equality, and
both composite isomorphisms have graph $x=y$. The two theories are
equiconsistent.
\end{theorem}
\begin{proof}
Compose the exact prefix/native correspondence with the pinned
native--$\Teq$ interpretation and the scalar--iPRA interpretation.
For completeness, the numerical round-trip argument is as follows.
External induction on each fixed scalar description shows that its
double translation satisfies its original recursion equations. Equation
induction, with the seed and passive parameters held fixed and only
afterward instantiated, proves uniqueness of two iterators with the same
seed and update. For each fixed
ordinary primitive-recursive description, external induction on that
description and internal uniqueness of primitive recursion give the
same conclusion. Hence each composite fixes every numerical operation
provably, at arbitrary internal arguments, and $x=y$ is the required
internally definable isomorphism. Formula round trips follow from the
native logical macro theorem and the inverse prefix translations.
Contradiction is preserved in both directions by (8) and
$\bot_*=0=S0$. Effective finite proof translations give equiconsistency.
This does not claim that either theory proves its own consistency.
\end{proof}

\section{Executable specification and measurement}
\label{sec:measurements}

\subsection{JSON syntax and binding discipline}

The standalone checker uses only the Python standard library.
Its numerical terms are encoded by
\begin{center}
\begin{tabular}{ll}
\toprule
Mathematical object & JSON encoding\\
\midrule
Zero & \code{0}\\
Bound variable, index $i>0$ & integer $i$\\
Free eigenparameter, level $i<0$ & integer $i$\\
Successor $St$ & \code{[t]}\\
Iterator $\I(n,a;u.h)$ & \code{[n,a,h]}\\
Formula $\Eq{k}{s}{t}$ & \code{[k,s,t]}\\
\bottomrule
\end{tabular}
\end{center}
The context determines whether a three-element array is a formula record
or an iterator term. The checker validates each use at its required sort.
Booleans and floating-point numbers are not numerical indices.

At outer scope $n$ and binder depth $d$, a numerical integer is admissible
exactly when $-n\le i\le d$. Only an iterator's update adds one to $d$.
A formula of count $k$ validates its leaf terms at depth $2k$.
The fresh parameter introduced by a local branch is $-n-1$.
All branches add exactly one context entry, so the implementation recovers
$n$ as the context length minus the initial entry.

Substitution inserts only terms that are closed with respect to exterior
positive indices. Such inserted terms may have their own internal
binders, but they contain no positive index referring to an ambient binder.
Consequently the insertion itself needs no extra lifting; the traversal
still increments the sought index under an iterator binder and decreases
larger indices when an external slot is removed. This invariant is central
to the short implementation and is checked at every certificate-supplied
witness. The homogeneous alternative additionally needs lifting when it
places an existing leaf term under a new iterator binder.

Proof nodes are
\begin{lstlisting}
["u", clause_index, target_term, minor_proof]
["v", major_equation_proof, minor_equation_proof]
["e", paired_formula_proof, branch_proof]
["i", paired_goal_formula, witness_term, branch_proof]
\end{lstlisting}
Input is exactly \code{\{"goal": F, "proof": P\}}. In a Use node whose
antecedent is reflexive at its leaf, \code{minor\_proof} is an inert JSON
field; the examples use \code{[]}. Its contents are not a claimed
subderivation. Every active premise, supplied term, introduced formula,
and elimination result is checked. The public goal must be a closed
well-scoped formula. The program prints \code{ACCEPT} on success and
\code{REJECT} on a detected exception; finite tests do not replace the
soundness proof. As with the comparison programs, completeness refers
to the ideal finite-derivation algorithm; a particular Python process
can exhaust its recursion, parsing, or memory limits on large inputs.
The binder count may encode a long finite prefix succinctly.

\subsection{Measured comparisons}

Measurements use CPython 3.12.14. Genuine module, function, and class
docstrings are removed, the complete module is normalized by
\code{ast.unparse} with one terminal newline, and significant Python
tokens, all AST nodes, and UTF-8 bytes are counted. Layout, comments,
encoding markers, and end markers are not significant tokens. Gzip level
9 with a fixed zero timestamp is informational. No gate, scope checker,
parser interface, or active inference routine is omitted from a measured
module.

\begin{center}
\small
\begin{tabular}{lrrrr}
\toprule
Complete implementation & Tokens & AST nodes & Bytes & Gzip-9\\
\midrule
Existing \code{uniform\_clauses} & 807 & 931 & 2912 & 994\\
Matched nested control & 791 & 890 & 2763 & 931\\
Prefix calculus $\F$ & \textbf{788} & 900 & \textbf{2700} & 941\\
Homogeneous prefix variant & 840 & 965 & 2843 & 979\\
\bottomrule
\end{tabular}
\end{center}

Relative to the exact existing four-family program, the delivered prefix
checker saves 19 tokens, 31 AST nodes, and 212 normalized bytes
(approximately 2.35\%, 3.33\%, and 7.28\%, respectively). This comparison
includes the simplified clause and certificate interfaces as well as the
prefix representation.

The matched nested control uses the same four proof tags, synthesized
elimination result, single-antecedent clauses, implicit scope protocol,
reflexive-leaf shortcut, and numerical term codec. It retains nested
Eq/$D$ formulas. Against this control, the prefix representation alone
saves 3 tokens and 63 bytes (approximately 2.28\% of bytes), while adding
10 AST nodes and 10 gzip bytes. Thus it is not a reduction in every
coordinate against that tighter control.

The antecedent-free implementation in the last row pays for its gate and
capture-free lifting. Its structural theorem does not yield the smallest
checker in this table. The counts also do not establish a new repository-wide
minimum: the existing validated native implementation has 788 \emph{bytes},
and a separate compact implementation has 516 bytes under its canonical
public-goal contract. Those frozen sources are included for transparency.
The number 788 in the prefix row is a \emph{token count}; the prefix
program has 2700 normalized bytes. These distinct units and contracts
must not be conflated.

The exact sources, source hashes, normalized copies, compressed copies,
and measurement procedure are in \code{code/} and
\code{code/evidence/checker\_metrics.json}.


\subsection{Verification and limits of the executable evidence}

The main reproduction run constructed 37 positive certificates and
translated each to the matched nested checker and the independently
retrieved existing four-family checker. Each ran in all three programs
under both ordinary Python and \code{python -O}. An additional 22 negative
inputs ran in both modes against the prefix checker. This is 266 subprocess
expectations, with zero mismatches. The positive cases include zero and
successor computation, derived symmetry, nested paired binders, and a
nontrivial equation-induction proof of
\[
 \I(y,0;u.Su)=y
\]
followed by specialization. The negative cases include incorrect induction
steps, escaped eigenparameters, malformed or out-of-scope terms, invalid
clause indices, nonintegral binder counts, and compound premises used as
equations.

There were also 2001 checks that opening commutes with expansion to the
nested representation. A separately written audit used environment-based
numerical semantics, rather than another syntactic substitution routine,
for 400 matrix openings, 400 witness openings, and 300 iterator-body
substitutions. It passed 20 additional rejection cases and 5 accepted
examples. The scope attacks include a forbidden parameter beneath a dead
zero-iterator body and beneath a formula binder; dead code is still
checked for scope.

The deterministic generators, concrete certificates, full results, and
source hashes are included. The independent semantic checks use bounded
samples and a finite test interpretation for formula-opening comparisons;
they are not asserted to be a model of the entire arithmetic theory.
Written derivation translations supply the metatheorems, while these
checks supply evidence that the implementations follow them.

From the package directory, run
\begin{lstlisting}
python3 code/reproduce_checker.py
python3 code/audit_flat.py code/flat_control.py
python3 code/flat_control.py < code/examples/equation_induction.json
\end{lstlisting}
The last command checks the supplied induction certificate. The example
filenames and complete format description are also listed in
\code{README.txt} and \code{code/FORMAT.txt}.


The correctness results in this report are written mathematical proofs.
The inherited native equivalences and arithmetic bi-interpretation have
Lean formalizations in the cited repository. This package does not claim
a newly checked Lean or Rocq formalization of its prefix checker or of
the homogeneous-context theorem. The Python programs verify certificates
for the presented inference rules; they are not proof assistants checking
the metatheory automatically.

\section{Further research questions}

\begin{enumerate}
\item Can the prefix representation reduce the smallest fully validated
compact checker while preserving that checker's exact public language and
successful-trace invariant? The matched readable reduction does not
settle this question.
\item Can homogeneous contexts be implemented without paying for an
additional gate operation, while keeping all compound assumptions and
their exact conditional meaning? The theorem in Appendix~\ref{sec:homogeneous}
shows that the antecedent field is mathematically eliminable, but the
measured implementation still pays more for that elimination.
\item Can a substitution-stable introduction, such as the existing
diagonal introduction, be paired with the prefix interface at lower total
checker cost than the witness-directed rule? Formula depth makes the
atomic test trivial, but it does not eliminate the arithmetic guard.
\item Can the prefix packing and opening lemmas, followed by the
homogeneous gate lemma, be added as short modules over the existing Lean
\code{CPrf}/\code{VPrf} development? This would check the new metatheory
without redoing the inherited arithmetic bridge.
\item What minimum is attainable under a fixed, explicit certificate
contract? A useful comparison should fix goal grammar validation, scope
checking, accepted padding, normalization, and the source measurement
procedure before optimizing implementations.
\end{enumerate}

\clearpage
\appendix
\section{The homogeneous four-family calculus}
\label{sec:homogeneous}

The proposed calculus retains the scalar terms and the original ordinary
meaning of the repository's Eq/$D$ language.  The change is in retained
assumptions: every context entry is a single one-variable formula
template.  No entry has a separately stored antecedent, a list of
antecedents, or a flag distinguishing facts from conditional closures.

\subsection{Language, ordinary reading, and the gate}

The term language is
\[
 t ::= x\mid 0\mid St\mid I(n,a;u.h),
\]
where $u$ is bound only in $h$.  Formulas are
\[
 A ::= (s=t)\mid D_{x,y}A.
\]
The ordinary reading is unchanged:
\[
 D_{x,y}B(x,y)\ :=\ \exists x\,
             \bigl(B(x,0)\to\forall y\,B(x,y)\bigr).
\]
Write $Q_yB$ as a metalinguistic abbreviation for
$B(0)\to\forall yB(y)$, not as a constructor of the new formula language.
A context $\Delta$ is a finite reusable list of templates $B(y)$; its
ordinary reading is the list of formulas $Q_yB$.  The distinguished
initial entry is
\[
 \rho(y)\equiv(y=y).
\]
It remains present in every context.  Its ordinary reading is a theorem.

For a fresh variable $v$, define
\[
 \kappa_v(s=t)
   :=\bigl(I(0,s;u.t)=I(v,s;u.t)\bigr),
 \qquad
 \kappa_v(D_{x,y}B):=D_{x,y}(\kappa_vB).
\]
Here $u$ is fresh and unused in $t$; all substitutions and insertions are
capture-free.  In a de Bruijn implementation, placing $t$ in the new
iterator body must lift its free indices.  The asymmetry is useful:
\[
 \kappa_0(s=t)\text{ is literally reflexive},\qquad
 \kappa_{Sz}(s=t)\ \dashv\vdash\ (s=t).
\]
The second statement uses iterator computation and identifies the
positive gate with the original equation.  Under nested
$D$ binders, the outside variable $v$ remains free; the gate does not
change the number of formula binders or an equation/compound root.

\subsection{The four rule families}

\paragraph{Use.}
For an entry $B(y)$ and a target term $t$, infer $B(t)$ from one minor
derivation of $B(0)$.  The minor may instead be omitted exactly when the
raw substituted base $B(0)$ has the form $D^j(r=r)$ with identical
leaf terms, for some $j\ge0$:
\[
 \frac{B(y)\in\Delta\qquad\Delta\vdash_H B(0)}
      {\Delta\vdash_H B(t)},
 \qquad
 \frac{B(y)\in\Delta\qquad B(0)\equiv D^j(r=r)}
      {\Delta\vdash_H B(t)}.
 \tag{Use}
\]
These are the one-premise and zero-premise cases of one instruction
family.  The omission test compares the two literal leaf terms and
does not inspect their numerical values.  It is valid at every binder
depth; an arbitrary compound theorem does not qualify.

\paragraph{Eval.}
From two equations, infer the result of one permitted computation:
\[
 \frac{\Delta\vdash_H s=t\qquad\Delta\vdash_H b=r}
      {\Delta\vdash_H r=\Phi(s,t,b)}.
 \tag{Eval}
\]
The partial operation $\Phi$ has exactly two cases:
\[
 \Phi(s,t,b)=
 \begin{cases}
 t,&b\equiv I(0,s;u.h),\\
 h[t/u],&b\equiv I(Sc,a;u.h)\ \text{and }s\equiv I(c,a;u.h).
 \end{cases}
\]
All other shapes are rejected.  In particular, no primitive root-joining
or arbitrary transport case is included.

\paragraph{Elim.}
\[
 \frac{\Delta\vdash_H D_{x,y}B(x,y)\qquad
       \Delta,B(b,y)\vdash_H A}
      {\Delta\vdash_H A}.
 \tag{Elim}
\]
The eigenparameter $b$ is fresh for the conclusion, original context and
fixed parameters.  The new retained entry is the single template
$y\mapsto B(b,y)$, with no separately stored zero instance.

\paragraph{Intro.}
Define
\[
 \eta(B,t,z)=
 \begin{cases}
 z,&B\text{ is a literal equation and }t\equiv0,\\
 0,&\text{otherwise}.
 \end{cases}
\]
Then
\[
 \frac{\Delta,\kappa_v(B(t,\eta(B,t,z)))\vdash_H B(t,Sz)}
      {\Delta\vdash_H D_{x,y}B(x,y)}.
 \tag{Intro}
\]
The eigenparameter $z$ is fresh for $\Delta,t$ and all fixed parameters
of $B$.  The displayed new entry is a template in its fresh argument $v$;
$z$ is a free parameter there.  The witness-directed choice of cases or
progress is exactly the repository's established witness rule.

There are no further primitive rules or nonlogical axioms.  In
particular, there is no Hyp rule and no Refl rule.  The designated
contradiction is $0=S0$.

\subsection{Recovering computation and numerical equality}

\paragraph{Reflexivity.}
Use the fixed entry $\rho(y)=(y=y)$ at target $t$.  Its base is literally
$0=0$, so no minor is required and the conclusion is $t=t$.


\paragraph{Reflexive prefixes.}
Every well-scoped $D^j(r=r)$ is derivable in every context.  Prove this
by induction on $j$, independently of the gate retrieval theorem.  The
case $j=0$ is the preceding derived reflexivity.  For $j>0$, apply Intro
with witness $1=S0$, which forces case mode.  Ignore its temporary gated
assumption.  The required successor target has $j-1$ binders, and its
leaf terms remain identical because both undergo the same substitutions.
The induction hypothesis supplies that target.  Consequently every
omitted Use minor can be expanded into an actual derivation using only
the fixed initial matrix and the displayed rules.

\paragraph{The two computations.}
Zero computation is Eval on the reflexivities $a=a$ and
$I(0,a;u.h)=I(0,a;u.h)$.  Successor computation is Eval on the
reflexivities of $I(c,a;u.h)$ and $I(Sc,a;u.h)$.

\paragraph{Root joining.}
From $s=t$ and $s=r$, put $Z=I(0,s;u.0)$.  Eval on $s=r$ and $Z=Z$
gives $Z=r$; Eval on $s=t$ and $Z=r$ gives $r=t$.  Symmetry is this
derived joining rule with one premise reflexive; transitivity follows by
one symmetry and one joining.

\paragraph{Transport in a term template.}
Suppose $s=t$ and $h(s)=r$.  Set
\[
 Z=I(0,s;u.h(u)),\qquad P=I(S0,s;u.h(u)).
\]
Zero-case Eval yields $Z=s$ and $Z=t$.  Successor-case Eval on $Z=s$
and $P=P$ gives $P=h(s)$, hence $P=r$ by the preceding derived equality
rules.  Successor-case Eval on $Z=t$ and $P=r$ now yields $r=h(t)$.
This is the native kernel's crossed one-template transport rule.
Consequently all numerical equality reasoning used below is derived
inside the four-family calculus.

\subsection{The gate retrieval theorem}
\label{subsec:gate-retrieval}

For every formula $A$ and every context containing $\rho$:
\begin{enumerate}
\item $\Delta\vdash_H\kappa_0A$.
\item A derivation of $\kappa_1A$ can be converted effectively to a
      derivation of $A$, and conversely.
\item If $\kappa_vA$ is a retained context entry, then
      $\Delta\vdash_H A$.
\end{enumerate}
The transformations are uniform under enlarging the context and under
renaming eigenparameters.

\paragraph{Proof of assertion 1.}
The formula $\kappa_0A$ has the same $D$ prefix as $A$ and literally
identical iterator terms at its leaf.  It is therefore an instance of
the independently established reflexive-prefix theorem.  No retrieval
lemma is used in this step.

\paragraph{Simultaneous proof of assertions 2 and 3.}
At an equation, assertion 2 is computation and equational transport.  Suppose assertion 3 has been proved for all formulas with
fewer binders than $A=D_{x,y}B$.

To convert $\kappa_1A=D_{x,y}\kappa_1B$ to $A$, eliminate the source
at a fresh $b$, retaining the matrix $\kappa_1B(b,y)$.  Introduce the
target with witness $b$.  Because $b$ is a variable rather than the
literal zero, this is a case introduction.  Its local entry is
$\kappa_v(B(b,0))$.  Assertion 3 at the smaller formula depth supplies
$B(b,0)$.  Assertion 2 at that smaller depth converts it to
$\kappa_1B(b,0)$.  Use the retained source matrix at target $Sz$ and
convert its result $\kappa_1B(b,Sz)$ back to $B(b,Sz)$.  Close Intro
and Elim.

The reverse conversion starts by eliminating $D_{x,y}B$ and introducing
$D_{x,y}\kappa_1B$ with the same fresh witness $b$.  Retrieve its local
zero assumption by assertion 3 at the smaller depth, convert to the
ungated zero instance, Use the retained $B(b,y)$ matrix at $Sz$, and
convert back to the gated successor instance.  Thus assertion 2 at the
current depth uses assertion 3 only at a strictly smaller depth.

Finally, prove assertion 3 at the current depth.  Use the retained entry
$\kappa_vA$ at target $1$.  Its zero instance is leaf-reflexive at every
formula depth, so the minor is omitted by the explicitly stated Use
rule.  The result is $\kappa_1A$, which assertion 2 converts to $A$.
This completes a well-founded simultaneous induction.  If desired, the
omitted minor can be expanded by assertion 1; no unproved compound
assumption has been discarded.

\subsection{Soundness in the original ordinary realization}

Read context entries as their conditional closures in $T_{\rm eq}$.
Formula induction in ordinary intuitionistic logic with equality proves
\[
 \kappa_0A,\qquad \kappa_{Sz}A\leftrightarrow A.
\]
At atoms the latter uses the two computations, and at a $D$ binder it
uses the ordinary extensionality of the displayed existential conditional
definition.  Numerical coverage therefore proves
\[
 A\ \longrightarrow\ Q_v\kappa_vA.
 \tag{Gate introduction}
\]
Indeed, under $A$, a numerical argument is zero or a successor; the
first gated instance is a theorem, and the second is equivalent to $A$.

Induct now on a four-family derivation.  Eval is a valid equational
consequence.  Use is conditional-universal specialization; when its
minor is omitted, the reflexive-prefix theorem supplies exactly that
missing minor.  Elim is ordinary existential elimination.  For Intro, its
ordinary temporary fact $B(t,\eta)$ supplies the retained encoded
assumption by Gate introduction.  When $\eta=0$, the successor premise
and numerical coverage establish the conditional universal conclusion.
When $\eta=z$, the matrix is an equation, and equation induction with
retained side assumptions establishes the same conditional conclusion.
Existential introduction at $t$ completes the argument.  The initial
$\rho$ closure is an equality theorem.  Hence the calculus is sound
under the unchanged ordinary reading.

\subsection{Completeness and the interpretation composites}

Let $W$ be the pinned seven-rule witness calculus \texttt{WPrf}.  The
repository proves \texttt{wprf\_iff\_nprf}; its introduction uses the
same witness test as above.  Translate its formulas by fixing equations
and $D$ and putting
\[
 (Q_yB)^\circ=D_{r,y}B^\circ
\]
with $r$ fresh and unused.  The translation preserves root class, fixed
parameters and capture-free numerical substitution.

The precise derivation induction is: if $\Gamma\vdash_W A$, and every
$C\in\Gamma$ has a proof of $C^\circ$ from a homogeneous context
$\Delta$, then $\Delta\vdash_H A^\circ$.  Hyp uses the supplied proof.
The computations and transport use the equational derivations above.
For $Q$ elimination, eliminate the translated dummy-witness $D$ and Use
its retained matrix at the original target.  For Intro, the gate
retrieval theorem supplies precisely the translated local fact from the
new entry $\kappa_v(B^\circ(t,\eta))$, so the induction hypothesis
applies to the source premise.

For source $D$ elimination, the newly retained matrix is
$B^\circ(b,y)$.  It supplies the translated source assumption
$(Q_yB(b,y))^\circ=D_{r,y}B^\circ(b,y)$ as follows: introduce this
dummy-witness binder at witness $1$, ensuring case mode, retrieve its
local zero fact through the gate theorem, and Use the retained matrix at
$Sz$.  Thus the derivation induction also applies to the source
elimination premise.  This strengthened context statement avoids any
unproved primitive cut or clause-construction rule.

For an explicit source context, take $\Delta$ to consist of $\rho$ and
the entries $\kappa_v(C^\circ)$ for $C\in\Gamma$; gate retrieval
provides all needed hypothesis proofs.  Hence every source derivation
has an effective four-family derivation.

Conversely, the soundness proof translates every new rule to the original
ordinary scalar theory.  The pinned native completeness theorem returns
it to the native calculus.  For context bookkeeping, the reading of a
matrix entry is exactly its native $Q$ formula, and
$Q_v\kappa_vA\leftrightarrow A$ follows from Gate introduction and
specialization at $1$.  Thus the encoded-context composites are also
provably equivalent, not merely equiconsistent.

On formulas, the $H$-to-native translation is the inclusion of the
Eq/$D$ language.  The other translation only replaces $Q$ by $D$ with
an unused witness.  The native theory proves this replacement
equivalent in both directions, so formula induction gives both
round trips.  The $H$-side formula composite is literally the identity.
Both term maps and both numerical domains are the identity; numerical
equality remains actual equality.  Composing with the repository's
native-to-iPRA bi-interpretation gives the requested bi-interpretability,
both numerical composite graphs $x=y$, and equiconsistency.  The
designated contradiction is fixed.  The arithmetic bridge and its
existing proof-assistant verification are inherited results; the new
gate and context theorem is a written proof unless separately formalized.

\subsection{A flat formula representation}

The Eq/$D$ grammar is unary above its atomic leaf.  Every formula is
uniquely of the form
\[
 D_{x_1,y_1}\cdots D_{x_k,y_k}(s=t).
\]
It may therefore be represented by $(k,s,t)$ instead of a recursively
nested formula tree, with the terms scoped under the $2k$ displayed
bound variables.  A gate leaves $k$ unchanged and transforms the single
leaf equation; a binder adds one to $k$.  Under a relative-index
representation, the gate's external variable must be placed outside all
$2k$ formula-bound variables, and placing the old right-hand term under the iterator binder
requires lifting its exterior positive indices. In the mixed convention
of the delivered checker, negative free-parameter levels remain unchanged.  This is a representation theorem with no change
of language or interpretation.  Checker measurements must still include
the integer and scope validation, capture-free substitutions and gate
construction.

\paragraph{Scope of the simplicity claim.}
The mathematical deletion is the separate antecedent field of every
retained clause.  Fact assumptions and conditional closures now use one
homogeneous context grammar.  The omitted-minor reflexivity test and the
gate are explicit parts of the trusted checker.  Four instruction
families do not beat the repository's existing four-family count by
themselves; any software-size improvement must be the result of a
matched, complete, measured implementation.


\paragraph{Implementation evidence.}
The separate homogeneous implementation was checked on 400 generated
formulas for its literally reflexive zero gate and on 800 positive-control
instances against environment-based numerical evaluation. Six closed
certificates passed, including equation induction whose local hypothesis
is recovered from its gated context entry, followed by specialization.
These finite checks are reproduced by \texttt{code/audit\_homogeneous.py};
they are separate from the primary checker's subprocess comparison.


\clearpage
\begin{thebibliography}{9}
\raggedright
\bibitem{repo} Vladimir Reshetnikov, \emph{Concise presentations of
intuitionistic primitive recursive arithmetic}, pinned repository README,
commit \code{17898ea7e6a1300dfdcad71040d5daeed67916d4}.
\url{https://github.com/VladimirReshetnikov/ProveIt/blob/17898ea7e6a1300dfdcad71040d5daeed67916d4/Logic/PrimitiveRecursiveArithmetic/ConciseIPRA/README.md}.
\bibitem{clausal} Same repository, \code{Lean/ConciseIPRA/NativeClausal.lean},
including \code{CPrf}, \code{cprf\_iff\_nprf}, and the clause translations.
\bibitem{uniform} Same repository,
\code{Lean/ConciseIPRA/NativeUniformClauses.lean} and
\code{Research/uniform\_clauses/code/checker.py}; the four-family clause
interface and its original standalone implementation.
\bibitem{eval} Same repository,
\code{Lean/ConciseIPRA/NativeEval.lean} and
\code{Research/clausal\_two\_case/}; computational equality and its
two-case reduction.
\bibitem{witness} Same repository,
\code{Lean/ConciseIPRA/NativeOneIntro.lean}; the exact equivalence of
the witness-directed single introduction and the native kernel.
\bibitem{native-bi} Same repository,
\code{Lean/ConciseIPRA/NativeBiInterpretation.lean},
\code{NativeRoundTrip.lean}, and \code{BiInterpretation.lean}; ordinary
readings, native round trips, and the scalar--iPRA bi-interpretation.
\bibitem{bridge} Same repository,
\code{Research/certified\_zero/reference/}: the inherited choice basis,
ordinary bridge, scalar predecessor, equality reflection, and primitive
recursion proofs. These five source files are included in this package
with provenance and the repository's MIT-0 license.
\end{thebibliography}
\end{document}
